\begin{abstract}
Crop rankings are ordinal, whereas acreage allocation is a cardinal decision made under joint uncertainty and operational constraints. We formalize this gap with a set-valued, multi-crop programme and distinguish possible, universal and selected ranking reversal. Exact results show that rankings alone do not identify acreage and that operational feasibility can force universal reversal even when downside risk is inactive. In pre-specified simulations, assigned operational interventions convert a corn-only allocation into soybean-majority allocations in every treated cell; all \EtwoPassedIntervals{}/\EtwoIntervals{} family-wise intervals meet the precision criterion. A separate information--flexibility experiment shows that their interaction can be positive (\EsixPositive{}, 95\% interval \EsixPositiveLow{} to \EsixPositiveHigh{}), exactly null or substitutive. Four other pre-specified experiments do not meet their experiment-level precision criteria and remain inconclusive. An official US panel extending from \EmpiricalStartYear{} to \EmpiricalEndYear{} contains \NStateYears{} complete state-years in \NStates{} states. Concurrent operating-margin inversion intensity is \OperatingInversion{} (state-cluster 95\% interval \OperatingInversionLow{} to \OperatingInversionHigh{}), but all four strictly lagged acreage-share intervals include zero and national summaries depend on aggregation. Ranking reversal is therefore demonstrable under explicit mechanisms, while public acreage disagreement alone does not identify private objectives, constraints, risk preferences, dependence, causality or welfare.
\end{abstract}
